\subsection{Central target extensions do not enlarge a cyclic-centralizer class action}

\begin{statement}{D23.P31}{Proved}{A limitation of central extensions of the cubic target}
\end{statement}
Let $G$ be a finite group, $g\in G$, and assume
$C_G(g)=\langle g\rangle$. For any central extension
\[
              1\longrightarrow Z\longrightarrow
              \widetilde G\xrightarrow{q}G\longrightarrow1
\]
and any lift $\widetilde g$ of $g$, its conjugacy class
$\widetilde C$ maps bijectively to the conjugacy class
$C$ of $g$. The bijection is equivariant for conjugation
and for every Hurwitz braid action on tuple powers.
Consequently, for every $n$, the braid actions on
$\widetilde C^n$ and $C^n$ are isomorphic, and so are
the actions after quotienting by simultaneous target
conjugation. Every individual braid orbit is preserved
with its exact marked permutation action.

In particular, replacing the cubic target
$E_0=\mathbb F_3^2\rtimes SL_2(\mathbb F_3)$ by any
central extension, while using one selected lift of
the original order-six conjugacy class, cannot enlarge
any of the Hurwitz permutation images considered here.
This includes a central extension of order $648$.

\begin{proof}
The inclusion
$C_{\widetilde G}(\widetilde g)\subseteq
q^{-1}(C_G(g))$ is immediate. Conversely, if
$q(h)=g^k$, then $h=\widetilde g^kz$ for some
$z\in Z$. The centrality of $Z$ implies that $h$
commutes with $\widetilde g$. Therefore
\[
 C_{\widetilde G}(\widetilde g)
                  =q^{-1}(\langle g\rangle).
\]
The induced map of coset spaces
\[
 \widetilde G/C_{\widetilde G}(\widetilde g)
                    \longrightarrow G/C_G(g)
\]
is a bijection, and these coset spaces parameterize
the respective conjugacy classes. This proves the
equivariant class bijection. Notice that no assumption
on the order of the chosen lift is needed.

Apply this bijection coordinatewise. The braid action
uses only $(u,v)\mapsto(uvu^{-1},u)$, which commutes
with the homomorphism $q$. Thus the product bijection
is braid equivariant. Central elements act trivially
by conjugation on $\widetilde C$, so its target action
factors through $G$ and agrees exactly with the action
on $C$. Passing to simultaneous conjugacy orbits
therefore remains a bijection. The same reasoning
shows that an individual braid orbit, including one
chosen by a fixed-product condition, has the same
marked permutation action as its image.

For the original cubic target, $|E_0|=216$, the chosen
element has order six, and its conjugacy class has
size $36$ by D23.P17's complete construction. Hence
its centralizer has order $216/36=6$ and contains
the cyclic subgroup of order six generated by it.
They are equal, so the general assertion applies.
\end{proof}

\paragraph{Exact limitation.}
This statement is about the specified selected
conjugacy-class Hurwitz construction. It does not say
that all evaluations in a central extension carry no
new information, nor does it cover noncentral target
extensions or a union of different lifted classes.
It rules out obtaining a new braid permutation image
merely by centrally extending this cyclic-centralizer
target and choosing one lifted class.
